\documentclass{article}
\usepackage{amsmath,hyperref}
\title{An exact old result at a late dock with two active byte windows}
\author{}
\date{}
\begin{document}
\maketitle
\paragraph{Audit.}
This development extension of the
\href{https://github.com/eternal-november/eos-research-programmes/pull/801}{late dock restart}
returns the actual historical result to Q, not merely a receipt that
some effect occurred. The station's abstract result is a fixture
calibration code in \(1,\ldots,256\), not a measured field observation.
It is stored at the real earlier native callback and is unavailable to
Q. SQLite represents the named digital effect. The selected stock
broker and receive bodies execute against the pinned Cave Circuit 01
geometry and radio configuration at the declared sites; this is not a
stock handoff protocol, full vehicle motion, power-loss test,
independent mission, or I1 acceptance.

\paragraph{What the handoff must distinguish.}
The same named report can be pending, or already credited with any of
256 known-domain results. Q's receipt tape and own debit are held
identical before the retry in those histories. C++ dispatch supplies an
actual prior result to a local cold-reopened journal if the earlier
report arrives. It is the value persisted under the same key that S
must return after a replay, not a recomputation or one fixed positive
symbol. The other owed inspection report still sends before
resolution. The protected logical dock maps to S after the late
departure of R and the native transfer callback, as before.

With no other readable side information, an exact handoff message must
distinguish 257 possibilities: the pending state and the 256 credited
values. An injective map into 256 octets is impossible. This is for
one nonempty late packet at a fixed cut, with known identity, key
and mode, and no independent timing, absence or sender-choice code.
It is not a limit on protocols that legitimately obtain another
observation. Saving and transporting prior returned results is already
part of
\href{https://web.stanford.edu/~ouster/cgi-bin/papers/rifl.pdf}{RIFL's completion-record and retry-rendezvous design}.
The added issue here is the exact incoming debit that migration
creates for the future receiver.

\paragraph{A qualified earlier byte.}
R can send S the high bit of its durable result code at \(-1.25\) on
the checked route. That earlier receive debit is no longer live in
S's one-second window at \(7/8\). Such a byte is not a completion
proof. At the actual later handoff, R reads the durable record
again, checks the earlier group still agrees, and emits a
single residual symbol. A designated zero represents pending;
each group has only 128 credited results. S decodes using the actual
earlier native callback, imports the result, then installs the single
logical binding; the retry at \(7/8\) yields the stored exact result
or commits one initial outcome. This uses an earlier byte \emph{and}
a later byte. It is a timing redistribution at the same total
unoptimized transfer length, not a free reduction in communication.

An actual later old packet, when admitted in a negative control,
can commit value 256 after an earlier prefix that described an
unpaid record. The checked interface refuses the stale combination;
using the old prefix without checking makes S return 128 instead
of 256 on a timely replay. The existing side report stays paid.
This is an example of invalidating the first byte's meaning, not loss
of radio receipt. A proof of a through-handoff update restriction
or a timely replacement prefix would instead be needed to promise
progress in that enlarged family.

\paragraph{Two separate native receive cuts.}
The three-byte retry and one-byte side send fit Q's cap four.
The two-byte direct handoff fills two bytes at S in the same window;
the retry is then charged at Q but cannot arrive at S, even
though the log and binding were installed correctly.
With a still-valid earlier group and one late residual byte,
\(1+3=4\) at S and the retry arrives. The explicit two-byte
record remains correct when the sender cap applied to that
receiver is five. A bit-only record can avoid a second effect yet
return the wrong result; the result is not equivalent to the
Boolean credit.

After S actually accepts the retry it sends the exact result back
at \(15/16\). This packet is charged against S's transmit budget and
Q's receive budget, under the pinned sender-cap rule. With three live
one-byte grants/preparation receipts at Q, an explicit two-byte
result fails against cap four even if S has already reconciled.
At this stage pending is no longer a possible result, so 256
exact values fit in one byte with the prebound reply format. The
short actual return reaches Q at the same cap. That format does not
claim to simultaneously report the first-versus-replay class or
recover arbitrary records. No caller confidence is inferred from
the local S result.

\paragraph{Checks and comparison.}
The sweep comprises 16,512 native traces. It checks all 257
prior states, both tested caps, both return widths, arrival of a
historical receipt, and the later wrong-prefix control.
In 514 comparable traces a one-byte reply arrives at the tight
return cut where a two-byte reply does not. The exact scoped
format handles all 257 fixed-scope states without raising the
tight retry budget. The full-record baseline has no such
retry at cap four; a direct policy using the same two-cut codec
and premises ties. The early pre-window path and the checked
prefix hold are comparative resources, so this does not prove an
advantage over all prior mechanisms. Lean verifies the finite
alphabet obstruction, the scoped codec, the stale-prefix result,
and the two byte frontiers. It does not verify C++, SQLite,
wall-clock process latency or motion. Run
\texttt{python3 result\_transfer.py} with the pinned parent inputs
or in the bundled offline replay.
\end{document}
